\chapter{The K3 Lattice, Periods and the Torelli Theorem}\label{ch:k3lattice}

\Cref{ch:k3} showed that all K3 surfaces are one and the same smooth four-manifold, with
$b_2=22$ and signature $-16$, while as complex manifolds they form a family with twenty
complex parameters. This chapter answers the question that these two facts raise:
\emph{how does one tell two K3 surfaces apart, and what does the space of all of them look
like?} The answer is a piece of linear algebra. The second cohomology $H^2(X,\Z)$ with its
intersection pairing is always the same lattice, $3U\oplus2E_8(-1)$ of \cref{ch:lattices}. A
complex structure is recorded by the position of one complex line, spanned by the class of
the holomorphic two-form, relative to that lattice; a Ricci-flat metric is recorded by the
position of a positive three-plane. The Torelli theorem says that nothing else is needed.

A student of string theory should care because this is the model for every moduli space that
follows. For a torus the moduli space was
$\Odd{d}\backslash\OddR{d}/(O(d)\times O(d))$: the positions of a positive $d$-plane relative
to the Narain lattice, up to lattice automorphisms (\cref{ch:narain,ch:odd}). For K3 we shall
find $O(\Gamma^{3,19})\backslash O(3,19)/(O(3)\times O(19))\times\R_+$, a space of the same
kind, and in \cref{ch:stringsk3} the $B$-field will enlarge $(3,19)$ to $(4,20)$. The points
of the moduli space at which the three-plane becomes orthogonal to a lattice vector of square
$-2$ are the K3 analogue of the self-dual radius of \cref{ch:tduality}: there the surface
develops an ADE singularity and the type~IIA string acquires a non-abelian gauge symmetry.

Conventions. $X$ is a complex K3 surface (\cref{ch:k3}), $\Lambda:=3U\oplus2E_8(-1)$ is the K3
lattice, also written $\Gamma^{3,19}$; $(x\cdot y)$ is the pairing, $x^2:=(x\cdot x)$, and
signatures are written $(n_+,n_-)$. The holomorphic two-form and its cohomology class are
both written $\Omega$ (Huybrechts writes $\sigma$). Vectors of positive square are called
space-like, following Aspinwall.

%=====================================================================================
\section{The intersection form on \texorpdfstring{$H^2(X,\Z)$}{H2(X,Z)}}
\label{sec:k3lattice:form}

\subsection{From cycles to a lattice}
On a compact oriented four-manifold two closed surfaces in general position meet in finitely
many points, and counting these points with orientation signs gives an integer
\begin{equation}\label{eq:k3lattice:intersection}
 \alpha_1\cdot\alpha_2=\#(\alpha_1\cap\alpha_2),\qquad \alpha_1,\alpha_2\in H_2(X,\Z)
 \qquad\text{\source{Asp §2.3, ll.~514--518}} .
\end{equation}
In terms of differential forms the same pairing is $(a\cdot b)=\int_X a\wedge b$ for closed
two-forms $a,b$; because two-forms commute under the wedge product the pairing is symmetric.
Since a K3 surface is simply connected, $H_2(X,\Z)$ has no torsion and is isomorphic to
$\Z^{22}$ as a group \source{Asp §2.3, ll.~513--514 and footnote~4}. A free abelian group
with a symmetric integer pairing is a lattice\index{lattice} in the sense of
\cref{ch:lattices}, provided the pairing is non-degenerate. We now determine its three
invariants: discriminant, parity and signature.

\paragraph{Unimodular, from Poincaré duality.}\index{Poincaré duality}
Poincaré duality states that for a basis $\{e_i\}$ of $H_2(X,\Z)$ there is a second basis
$\{e_j^*\}$ with
\begin{equation}\label{eq:k3lattice:PD}
 e_i\cdot e_j^*=\delta_{ij}\qquad\text{\source{Asp §2.3, ll.~536--545}} .
\end{equation}
Let $G_{ij}=e_i\cdot e_j$ be the Gram matrix and write $e_j^*=\sum_kP_{kj}e_k$ with an integer
matrix $P$. Then \eqref{eq:k3lattice:PD} reads $GP=\mathbf 1$: the Gram matrix has an integer
inverse. By \cref{ch:lattices} this is one of the equivalent forms of
unimodularity\index{unimodular lattice}: $\det G\cdot\det P=1$ in $\Z$ forces $\det G=\pm1$,
and the dual lattice $\Lambda^*$ equals $\Lambda$. Duality also identifies the lattice of
integral homology $H_2(X,\Z)$ with that of integral cohomology $H^2(X,\Z)$, and from now on
we do not distinguish them. Huybrechts quotes the same fact in the form: the intersection
form of a compact oriented four-manifold, modulo torsion, is unimodular
\source{Huy ch.~1 §3.3, ll.~586--589}.

\paragraph{Even, from Wu's formula.}\index{Wu's formula}\index{even lattice}
The intersection form of a compact smooth four-manifold $M$ is even if and only if the second
Stiefel--Whitney class $w_2(M)$ vanishes, and for any almost complex structure
$w_2(M)\equiv c_1\pmod 2$ \source{Huy ch.~1, proof of Prop.~3.5, ll.~610--613}. The mechanism
behind the first statement is the congruence $x^2\equiv(x\cdot w_2)\pmod2$ for all
$x\in H^2(M,\Z)$ (this is the usual form of Wu's formula; it is standard, and our sources
quote only its consequence). For a K3 surface $c_1=0$, so
\begin{equation}\label{eq:k3lattice:even}
 x^2\in2\Z\qquad\text{for all }x\in H^2(X,\Z)\qquad\text{\source{Asp §2.3, ll.~546--552}} .
\end{equation}
The condition $w_2=0$ is the condition for $M$ to admit spinors. The evenness of the K3
lattice is therefore the lattice-theoretic trace of the fact that fermions can be defined on
a K3 surface.

A student can test \eqref{eq:k3lattice:even} on holomorphic curves without knowing Wu's
formula. For a smooth curve $C$ of genus $g$ in a surface with trivial canonical bundle the
adjunction formula gives
\begin{equation}\label{eq:k3lattice:adjunction}
 C^2=2g-2\qquad\text{\source{Asp §2.6, eq.~(53), ll.~1178--1184}} ,
\end{equation}
because the degrees of the normal bundle and of the tangent bundle of $C$, which are $C^2$
and $2-2g$, must add up to $c_1(X)|_C=0$. On a quartic $X\subset\mathbb{P}^3$ a plane section
is a smooth plane quartic curve, of genus $(4-1)(4-2)/2=3$, so $C^2=4$; and indeed two planes
meet in a line, which meets the quartic in four points \source{Asp §2.5, ll.~964--967}. A
sphere ($g=0$) always has $C^2=-2$. Such curves will be the main actors of
\cref{sec:k3lattice:roots}. This test is not a proof of evenness, because for a fixed complex
structure the classes of curves do not span $H^2(X,\Z)$
\source{Huy ch.~1, footnote~11, ll.~641--643}.

\paragraph{Signature $(3,19)$, by two routes.}\index{signature}
The signature theorem gives $b_+-b_-=\tfrac13(c_1^2-2c_2)=-\tfrac23\chi(X)=-16$
\source{Asp §2.3, eq.~(20), ll.~527--535; Huy ll.~614--618}, as computed in \cref{ch:k3}. With
$b_++b_-=b_2=22$ one finds $(b_+,b_-)=(3,19)$. The second route uses Hodge theory and tells
us \emph{which} three directions are positive. Decompose
\begin{equation}\label{eq:k3lattice:hodgesplit}
 H^2(X,\R)=\bigl(H^{2,0}\oplus H^{0,2}\bigr)_\R\;\oplus\;H^{1,1}_\R,\qquad
 \dim=2h^{2,0}+h^{1,1}=2+20 .
\end{equation}
On the first summand, spanned by $\operatorname{Re}\Omega$ and $\operatorname{Im}\Omega$, the
form is positive definite (we prove this in \cref{sec:k3lattice:periods}). On $H^{1,1}_\R$ the
Hodge index theorem allows exactly one positive direction, that of a Kähler class
\source{Huy ch.~1, Remark~3.4, ll.~553--557}. Hence
\begin{equation}\label{eq:k3lattice:bpm}
 b_+=2h^{2,0}+1=3,\qquad b_-=h^{1,1}-1=19 .
\end{equation}
The two routes use different inputs (characteristic classes; Hodge numbers) and agree. The
second one is what the Lean theorem \lean{sigK3\_matches\_hodge} of
\cref{sec:k3lattice:lean} compares with the lattice decomposition.

\subsection{The theorem and what it does not say}
An indefinite even unimodular lattice is determined up to isomorphism by its signature, and
exists only if $n_+-n_-\equiv0\pmod8$ (\cref{ch:lattices};
\source{Asp §2.3, eq.~(23), ll.~557--566}). For $(3,19)$ we have $3-19=-16$, and the lattice
with that signature built from the two blocks of \cref{ch:lattices} is
$U^{\oplus3}\oplus E_8(-1)^{\oplus2}$, with $3\cdot(1,1)+2\cdot(0,8)=(3,19)$.

\begin{theorem}[The K3 lattice; Tier~L]\label{thm:k3lattice:lattice}\index{K3 lattice}
For every complex K3 surface $X$ there is an isometry of lattices
\begin{equation}\label{eq:k3lattice:lattice}
 H^2(X,\Z)\;\cong\;\Lambda:=E_8(-1)\oplus E_8(-1)\oplus U\oplus U\oplus U .
\end{equation}
\TL{} \source{Huy ch.~1, Prop.~3.5, eq.~(3.4), ll.~600--618; Asp §2.3, eq.~(24), ll.~567--631}
\end{theorem}

\begin{example}[The Gram matrix, checked by machine arithmetic]\label{ex:k3lattice:gram}
In a basis adapted to \eqref{eq:k3lattice:lattice} the Gram matrix is the block-diagonal
$22\times22$ integer matrix $G_\Lambda=\diag(-E_8,-E_8,U,U,U)$, where $E_8$ is the Cartan
matrix of \cref{ch:lattices} and
$U=\bigl(\begin{smallmatrix}0&1\\1&0\end{smallmatrix}\bigr)$. Its invariants follow block by
block: the diagonal entries are $-2$ and $0$, hence even; $\det G_\Lambda=(\det(-E_8))^2(\det
U)^3=(+1)^2(-1)^3=-1$; each $U$ has eigenvalues $\pm1$ and each $-E_8$ is negative definite,
giving $3$ positive and $16+3=19$ negative eigenvalues. A symbolic check (sympy, exact
determinant; numerical eigenvalues) confirms $\det=-1$ and the inertia $(3,19)$. In Lean the
blocks are certified separately (\cref{ch:lattices}) and the signature arithmetic is the
theorem \lean{sigK3\_eq}.
\end{example}

\begin{pitfallbox}[``Abstractly isomorphic'' is not ``equal'']
\Cref{thm:k3lattice:lattice} does not hand us $22$ surfaces inside $X$ that intersect
according to $G_\Lambda$. It says that such a basis of \emph{classes} exists, and the proof
finds it by classification, not by geometry. Huybrechts remarks that it would be interesting
to see the decomposition directly on some K3 surface by writing down cycles, and proposes
the Kummer surface as a candidate \source{Huy ll.~620--623}; see \cref{ch:kummer}. Moreover,
the isometry in \eqref{eq:k3lattice:lattice} is far from unique: it can be composed with any
element of the infinite group $O(\Lambda)$. Choosing one is extra data, called a marking
(\cref{def:k3lattice:marking}), and keeping track of this choice is the main bookkeeping
task of the chapter.
\end{pitfallbox}

A second caveat concerns the converse. A simply connected complex surface whose
intersection form is \eqref{eq:k3lattice:lattice} need not be diffeomorphic to a K3 surface,
although by Freedman's results it is homeomorphic to one
\source{Huy ch.~1, Remark~3.6(ii), ll.~633--639}. The lattice determines the topological type
of the manifold, but not its smooth structure.

%=====================================================================================
\section{Periods}\label{sec:k3lattice:periods}

\subsection{The two relations satisfied by \texorpdfstring{$\Omega$}{Omega}}
The lattice is the same for all K3 surfaces, so it cannot distinguish them. What varies is
the position of the Hodge decomposition \eqref{eq:k3lattice:hodgesplit} relative to the
lattice. To measure it, choose a basis $\{e_i\}$ of two-cycles and integrate the holomorphic
two-form:
\begin{equation}\label{eq:k3lattice:periods}
 \varpi_i=\int_{e_i}\Omega ,\qquad i=1,\dots,22
 \qquad\text{\source{Asp §2.3, eq.~(26), ll.~651--655}} .
\end{equation}
These $22$ complex numbers are the \emph{periods}\index{period} of $X$. They are the
coordinates of the class $\Omega\in H^2(X,\C)=\Lambda\otimes\C$ in the basis dual to
$\{e_i\}$. Since $\Omega$ is defined only up to a nonzero complex factor (\cref{ch:k3}), the
meaningful object is the point $[\Omega]=[\varpi_1:\dots:\varpi_{22}]$ of the projective
space $\mathbb{P}(\Lambda_\C)\cong\mathbb{P}^{21}$.

Not every point of $\mathbb{P}^{21}$ can occur. Write $\Omega=x+\ii y$ with
$x,y\in H^2(X,\R)$. In local coordinates $\Omega=f\,\dd z_1\wedge\dd z_2$, so
$\Omega\wedge\Omega$ is a $(4,0)$-form and vanishes on a complex surface, where only
$(2,2)$-forms can be integrated. Hence
\begin{equation}\label{eq:k3lattice:OmegaOmega}
 0=\int_X\Omega\wedge\Omega=(x+\ii y)\cdot(x+\ii y)=(x^2-y^2)+2\ii\,(x\cdot y)
 \qquad\text{\source{Asp §2.3, eq.~(29), ll.~672--682}} ,
\end{equation}
so $x^2=y^2$ and $x\cdot y=0$. On the other hand $\Omega\wedge\bar\Omega
=|f|^2\,\dd z_1\wedge\dd z_2\wedge\dd\bar z_1\wedge\dd\bar z_2$ is a positive multiple of the
volume form (in real coordinates $z_k=u_k+\ii v_k$ it equals
$4|f|^2\,\dd u_1\wedge\dd v_1\wedge\dd u_2\wedge\dd v_2$), so
\begin{equation}\label{eq:k3lattice:OmegaOmegabar}
 0<\int_X\Omega\wedge\bar\Omega=(x+\ii y)\cdot(x-\ii y)=x^2+y^2
 \qquad\text{\source{Asp §2.3, eq.~(30), ll.~683--694}} .
\end{equation}
Together: $x^2=y^2>0$ and $x\cdot y=0$. The real and imaginary parts of $\Omega$ are
orthogonal vectors of equal positive length. They span a plane in $H^2(X,\R)\cong\R^{3,19}$
on which the form is positive definite, a \emph{space-like two-plane}, and the ordered pair
$(x,y)$ orients it. Rescaling $\Omega\mapsto\lambda\Omega$ with $\lambda=r\ee^{\ii\theta}$
rotates $(x,y)$ by the angle $\theta$ and stretches it by $r$: plane and orientation are
unchanged. Complex conjugation $(x,y)\mapsto(x,-y)$ reverses the orientation
\source{Asp §2.3, ll.~697--702}. This proves the claim used for
\eqref{eq:k3lattice:bpm}, and it shows that the data carried by the periods is exactly an
oriented positive two-plane.

\subsection{The period domain}
\begin{definition}[Period domain]\label{def:k3lattice:D}\index{period domain}
For a lattice $\Lambda$ of signature $(n_+,n_-)$ with $n_+\ge2$,
\begin{equation}\label{eq:k3lattice:D}
 D:=\bigl\{[x]\in\mathbb{P}(\Lambda_\C)\;\big|\;x^2=0,\ (x\cdot\bar x)>0\bigr\}
 \qquad\text{\source{Huy ch.~6 §1.1, ll.~4281--4290}} .
\end{equation}
\end{definition}
The second condition is well posed because $(\lambda x\cdot\overline{\lambda x})
=|\lambda|^2(x\cdot\bar x)$. For the K3 lattice, $D$ is an open subset (in the classical
topology) of a smooth quadric hypersurface in $\mathbb{P}^{21}$, hence a complex manifold of
dimension $21-1=20$. This is the number $h^{1,1}=20$ of complex-structure deformations found
in \cref{ch:k3}, and it is the first sign that periods are good coordinates.

\begin{proposition}[Three descriptions of $D$; Tier~L]\label{prop:k3lattice:grass}
The map $[x]\mapsto\R\operatorname{Re}x\oplus\R\operatorname{Im}x$ identifies $D$ with the
Grassmannian $\mathrm{Gr}^{po}(2,\Lambda_\R)$ of oriented positive two-planes, and
\begin{equation}\label{eq:k3lattice:homog}
 D\;\cong\;\mathrm{Gr}^{po}(2,\Lambda_\R)\;\cong\;O(n_+,n_-)\big/\bigl(SO(2)\times
 O(n_+-2,n_-)\bigr) .
\end{equation}
\TL{} \source{Huy ch.~6, Prop.~1.5, ll.~4347--4363; Asp §2.3, eq.~(31), ll.~726--747}
\end{proposition}
\begin{proof}[Idea]
We have already seen that $x^2=0$, $(x\cdot\bar x)>0$ make $\operatorname{Re}x$,
$\operatorname{Im}x$ an orthogonal basis of equal positive length of a plane, and that
$\lambda x$ gives the same oriented plane. Conversely an oriented positive plane with oriented
orthonormal basis $e_1,e_2$ is the image of $e_1+\ii e_2$. For the second isomorphism,
$O(n_+,n_-)$ acts transitively on oriented positive planes, and the stabilizer of one of them
consists of rotations inside the plane, $SO(2)$, times arbitrary isometries of its orthogonal
complement, which has signature $(n_+-2,n_-)$.
\end{proof}

As a check, count real dimensions with $\dim O(p,q)=\tfrac12(p+q)(p+q-1)$:
\begin{equation}\label{eq:k3lattice:dim40}
 \dim_\R D=\dim O(3,19)-\dim SO(2)-\dim O(1,19)=231-1-190=40=2\cdot20 .
\end{equation}
The count can also be done by hand: a two-plane in $\R^{22}$ is given by two vectors ($44$
numbers) up to $GL(2,\R)$ ($4$ numbers), and positivity is an open condition.

A point of $D$ determines the whole Hodge decomposition\index{Hodge structure of K3 type},
not only $H^{2,0}$: set $H^{2,0}=\C x$, $H^{0,2}=\C\bar x$, and let $H^{1,1}$ be the
orthogonal complement of both. This is the unique Hodge structure on $\Lambda$ with
$H^{2,0}=\C x$ for which $H^{1,1}\perp H^{2,0}$ \source{Huy ch.~6, Prop.~1.2, ll.~4304--4318}.
The orthogonality is geometric: for a $(1,1)$-form $a$ the product $a\wedge\Omega$ is of type
$(3,1)$ and vanishes. In particular the period determines which integral classes are of type
$(1,1)$:
\begin{equation}\label{eq:k3lattice:NS}
 \mathrm{NS}(X)=\mathrm{Pic}(X)=H^2(X,\Z)\cap H^{1,1}(X)
 =\{v\in H^2(X,\Z)\mid(v\cdot\Omega)=0\} ,
\end{equation}
the \emph{Picard} or Néron--Severi lattice\index{Picard lattice}\index{Néron--Severi
lattice}, of rank $\rho(X)\le20$ \source{Huy ch.~1, eqs.~(3.2)--(3.3), ll.~570--575; Asp §2.5,
eq.~(41), ll.~921--930}. (The second equality in \eqref{eq:k3lattice:NS} holds because an
integral class is real, so $(v\cdot\Omega)=0$ implies $(v\cdot\bar\Omega)=0$.) The condition
$(v\cdot\Omega)=0$ is one complex linear equation on the period for each $v$. A generic
period satisfies none of them, and the corresponding K3 surface contains no curves at all:
$\rho=0$ \source{Asp §2.5, ll.~931--935}. If $X$ is algebraic, $\mathrm{NS}(X)$ has signature
$(1,\rho-1)$, because it is orthogonal to the positive plane of $\Omega$ and contains the
class of a hyperplane section \source{Asp §2.5, ll.~1018--1022}. Its orthogonal complement
in $H^2(X,\Z)$ is the \emph{transcendental lattice}\index{transcendental lattice} $T(X)$, of
signature $(2,20-\rho)$; it is rarely unimodular \source{Asp §2.5, ll.~1027--1029}.

\begin{example}[A period point with $\rho=20$]\label{ex:k3lattice:rho20}
Let $e_k,f_k$ ($k=1,2,3$) be the standard bases of the three copies of $U$:
$e_k^2=f_k^2=0$, $e_k\cdot f_k=1$. Put
\[
 x=e_1+f_1,\qquad y=e_2+f_2,\qquad \Omega=x+\ii y .
\]
Then $x^2=2(e_1\cdot f_1)=2$, $y^2=2$ and $x\cdot y=0$ because the two copies of $U$ are
orthogonal. Therefore $\Omega^2=x^2-y^2+2\ii\,x\cdot y=0$ and
$\Omega\cdot\bar\Omega=x^2+y^2=4>0$: the point $[\Omega]$ lies in $D$. By
\eqref{eq:k3lattice:NS} the Picard lattice consists of the integral vectors orthogonal to $x$
and $y$. A vector $ae_1+bf_1$ has $(ae_1+bf_1)\cdot x=a+b$, so the orthogonal complement of
$x$ in the first $U$ is spanned by $e_1-f_1$, of square $-2$. Hence
\[
 \mathrm{NS}=\Z(e_1-f_1)\oplus\Z(e_2-f_2)\oplus U\oplus E_8(-1)^{\oplus2},\qquad
 \rho=1+1+2+16=20,
\]
of signature $(1,19)$ as predicted, and $T=\Z x\oplus\Z y$ with Gram matrix
$\bigl(\begin{smallmatrix}2&0\\0&2\end{smallmatrix}\bigr)$, of determinant $4\ne\pm1$. The
period has the largest possible Picard number. That a K3 surface with this period exists is
guaranteed by \cref{thm:k3lattice:surj} below; we did not have to write down an equation. We
return to this surface in \cref{ex:k3lattice:maximal}. (Symbolic check of all pairings:
sympy.)
\end{example}

\begin{remark}[The case $n_+=2$: upper half-planes]\label{rem:k3lattice:tube}
For algebraic K3 surfaces the hyperplane class $\ell$ stays of type $(1,1)$, the period is
confined to $\ell^\perp$ of signature $(2,19)$, and the quartic surfaces, for example, have
the moduli space $O(\Lambda_C)\backslash O(2,19)/(O(2)\times O(19))$ of complex dimension
$\tfrac12(210-1-171)=19$. This matches the $35-16=19$ coefficients of a quartic polynomial
modulo $GL(4,\C)$; one of the $20$ deformations is lost
\source{Asp §2.5, eq.~(44), ll.~968--983}. For $n_+=2$ the period domain has a familiar
shape. Split $\Lambda_\R=U_\R\oplus W$ and send $\zeta\in W_\C$ to
$x=e-\tfrac12\zeta^2f+\zeta$. Then $x^2=2\cdot(-\tfrac12\zeta^2)+\zeta^2=0$ identically, and
with $\zeta=a+\ii b$ one finds $(x\cdot\bar x)=-\operatorname{Re}\zeta^2+\zeta\cdot\bar\zeta
=-(a^2-b^2)+(a^2+b^2)=2b^2$. So $D$ is the tube domain $\{\zeta\mid(\operatorname{Im}
\zeta)^2>0\}$ (this is \source{Huy ch.~6, Prop.~1.7, ll.~4384--4396} with $\zeta=\sqrt2z$).
For $W$ of signature $(1,0)$ this is the union of the upper and the lower half-plane: the
parameter $\tau$ of \cref{ch:torus2} is a period in exactly this sense, and $D$ for
signature $(2,n_-)$ is its higher-dimensional analogue. For $n_+=2$, $D$ has two
components exchanged by complex conjugation; for the full K3 lattice, $n_+=3$, it is
connected \source{Huy ch.~6, Remark~1.6, ll.~4365--4373}.
\end{remark}

%=====================================================================================
\section{Markings, the period map and the Torelli theorems}\label{sec:k3lattice:torelli}

\subsection{Marked K3 surfaces}
To compare the periods of two different surfaces they must be written in the same
coordinates, that is, in the fixed lattice $\Lambda$.

\begin{definition}[Marking]\label{def:k3lattice:marking}\index{marking}\index{marked K3 surface}
A \emph{marking} of a K3 surface $X$ is an isometry $\varphi\colon H^2(X,\Z)\to\Lambda$. Two
marked surfaces $(X,\varphi)$, $(X',\varphi')$ are equivalent if there is a biholomorphic
map $g\colon X\to X'$ with $\varphi\circ g^*=\varphi'$. The set of equivalence classes is
denoted $N$. \source{Huy ch.~7 §2.1, ll.~5600--5606; Asp §2.3, ll.~656--660}
\end{definition}
Concretely, a marking is a choice of which $22$ cycles are used in
\eqref{eq:k3lattice:periods}. The \emph{period map}\index{period map} is
\begin{equation}\label{eq:k3lattice:periodmap}
 \mathcal P\colon N\to D\subset\mathbb P(\Lambda_\C),\qquad
 (X,\varphi)\mapsto[\varphi_\C(\Omega_X)] ,
\end{equation}
and by \cref{sec:k3lattice:periods} it does take values in $D$. Three theorems describe it.
All three are Tier~L; none is formalized in this book's Lean libraries, and none could be
stated there at present, because the libraries contain no complex manifolds.

\begin{theorem}[Local Torelli; Tier~L]\label{thm:k3lattice:local}\index{Torelli theorem!local}
A K3 surface $X_0$ has a smooth universal deformation space $\mathrm{Def}(X_0)$ of complex
dimension $20$, and the period map $\mathcal P\colon\mathrm{Def}(X_0)\to D$ is a local
isomorphism. \TL{} \source{Huy ch.~6, Cor.~2.7 and Prop.~2.8, ll.~4646--4670}
\end{theorem}
The proof explains why it is true. First-order deformations of the complex structure are
classified by $H^1(X_0,T_{X_0})$, where $T$ is the holomorphic tangent bundle. The tangent
space of $D$ at the line $\ell=H^{2,0}$ is $\mathrm{Hom}(\ell,\ell^\perp/\ell)\cong
\mathrm{Hom}(H^{2,0},H^{1,1})$ \source{Huy ch.~6, Remark~1.1, ll.~4296--4302}, and the
differential of $\mathcal P$ sends a deformation $v$ to the map $\Omega\mapsto v\lrcorner
\Omega$. Because $\Omega$ vanishes nowhere, contraction with it is an isomorphism of bundles
$T_X\to\Omega^1_X$, hence an isomorphism $H^1(T_X)\to H^1(\Omega^1_X)=H^{1,1}$ of
$20$-dimensional spaces. In words: to first order a change of complex structure is the same
thing as a change of the class of $\Omega$ in the $(1,1)$ directions. Gluing the local
deformation spaces makes $N$ a $20$-dimensional complex manifold, which however is not
Hausdorff \source{Huy ch.~7 §2.1, ll.~5607--5616}.

\begin{theorem}[Surjectivity of the period map; Tier~L]\label{thm:k3lattice:surj}
\index{period map!surjectivity}
On each connected component $N^o$ of $N$ the period map $\mathcal P\colon N^o\to D$ is
surjective: every point of the period domain is the period of a marked K3 surface. \TL{}
\source{Huy ch.~7, Thm.~4.1, ll.~5892--5895; first proved by Todorov, ll.~5884--5885}
\end{theorem}

\begin{theorem}[Global Torelli; Tier~L]\label{thm:k3lattice:torelli}
\index{Torelli theorem!global}\index{Hodge isometry}
Two complex K3 surfaces $X,X'$ are isomorphic if and only if there is a \emph{Hodge isometry}
$\psi\colon H^2(X,\Z)\to H^2(X',\Z)$, that is, an isometry of lattices whose complexification
maps $H^{2,0}(X)$ to $H^{2,0}(X')$. Moreover, if $\psi$ maps at least one Kähler class of $X$
to a Kähler class of $X'$, then $\psi=f^*$ for a unique isomorphism $f\colon X'\to X$. \TL{}
\source{Huy ch.~7, Thm.~5.3, ll.~5969--5972}
\end{theorem}
In the language of markings: if $\mathcal P(X,\varphi)=\mathcal P(X',\varphi')$ then
$\psi=\varphi'^{-1}\circ\varphi$ is a Hodge isometry, so $X\cong X'$. The period is a complete
invariant of the complex structure. Aspinwall states the combination of
\cref{thm:k3lattice:surj,thm:k3lattice:torelli} as one sentence: the moduli space of complex
structures on a marked K3 surface, including orbifold points, is the space of possible
periods \source{Asp §2.3, Theorem~2, ll.~721--725}.

\begin{pitfallbox}[A Hodge isometry need not come from a map of surfaces]
The first half of \cref{thm:k3lattice:torelli} says that $X$ and $X'$ are isomorphic; it does
not say that the given $\psi$ is induced by an isomorphism. The isometry $-\mathrm{id}$ is a
Hodge isometry of every K3 surface and acts trivially on $D$, but it is never of the form
$f^*$, because it maps Kähler classes to their negatives
\source{Huy ch.~7 §5.2, ll.~5982--5984}. The reflections of \cref{sec:k3lattice:roots} in
classes of curves are further examples. This is the reason for the Kähler-cone condition in
the second half, and the reason why $N$ fails to be Hausdorff: several marked surfaces,
differing by such isometries, sit over the same period point and cannot be separated.
\end{pitfallbox}

\subsection{Forgetting the marking}
To obtain the moduli space of unmarked surfaces one divides by the changes of marking. Which
elements of $O(\Lambda)$ should be regarded as trivial? Those induced by diffeomorphisms of
the underlying four-manifold. Let $O^+(\Lambda)\subset O(\Lambda)$ be the subgroup of index
two that preserves the orientation of the positive three-planes; equivalently, the
isometries of spinor norm $+1$ \source{Huy ch.~7 §5.4, ll.~6031--6041}. Every element of
$O^+(\Lambda)$ is induced by a diffeomorphism and no element outside it is
\source{Asp §2.3, ll.~757--761}; Huybrechts proves the closely related statement that the
monodromy group of K3 surfaces is exactly $O^+(H^2(X,\Z))$
\source{Huy ch.~7, Prop.~5.5, ll.~6049--6051}. Therefore the moduli space of complex
structures on a K3 surface, including orbifold points, is
\begin{equation}\label{eq:k3lattice:Mc}
 \mathcal M_c\;\cong\;O^+(\Gamma^{3,19})\big\backslash O^+(3,19)\big/
 \bigl(O(2)\times O(1,19)\bigr)^+
 \qquad\text{\TL\ \source{Asp §2.3, eq.~(32), ll.~762--766}} ,
\end{equation}
where the superscripts $+$ restrict to the index-two subgroups preserving the orientation of
the space-like directions. The structure is that of \cref{ch:odd}: a Teichm\"uller space
$D$ (all periods), divided by a modular group $O^+(\Gamma^{3,19})$ (relabellings of the
cycles). But there is a difference of substance. The action of $O^+(\Gamma^{3,19})$ on $D$ is
ergodic, and $\mathcal M_c$ is not Hausdorff \source{Asp §2.3, ll.~767--769; Huy ch.~6 §1.3,
ll.~4409--4412}. Underlying this is the fact that the stabilizer $SO(2)\times O(1,19)$ of a point
is not compact, so that the discrete group need not act properly discontinuously. A
geometric symptom: a K3 surface may have an infinite automorphism group, which then fixes
its period point \source{Huy ch.~6, Example~1.9, ll.~4414--4420}. Aspinwall notes that
$\mathcal M_c$ does not seem to make a natural appearance in string theory, whereas the
moduli spaces which do appear are Hausdorff \source{Asp §2.3, ll.~770--774}. The next two
sections construct the first of them.

%=====================================================================================
\section{\texorpdfstring{$(-2)$}{(-2)}-curves, reflections and the Kähler cone}
\label{sec:k3lattice:roots}

A complex structure fixes the plane of $\Omega$. A Kähler metric adds a Kähler class
$J\in H^{1,1}_\R$ with $J^2=\int_XJ\wedge J>0$. Which classes are Kähler? By
\eqref{eq:k3lattice:bpm} the form on $H^{1,1}_\R\cong\R^{1,19}$ is Lorentzian, so the set
$\{\alpha\mid\alpha^2>0\}$ is the interior of a light-cone with two components. The
\emph{positive cone}\index{positive cone} $\mathcal C_X$ is the component containing a Kähler
class, and the Kähler cone\index{Kähler cone} satisfies $\mathcal K_X\subset\mathcal C_X$
\source{Huy ch.~7 §3.2, ll.~5812--5818}. It is cut out of $\mathcal C_X$ by the spheres
contained in $X$.

\begin{definition}[Roots]\label{def:k3lattice:roots}\index{root}\index{$(-2)$-curve}
A \emph{root} of a lattice is a vector $\delta$ with $\delta^2=-2$. A
\emph{$(-2)$-curve} on $X$ is a smooth rational curve $C\cong\mathbb P^1\subset X$; by
\eqref{eq:k3lattice:adjunction} its class is a root of $\mathrm{NS}(X)$.
\end{definition}
There is a converse: if $\delta\in\mathrm{NS}(X)$ is a root, then $\delta$ or $-\delta$ is
the class of an effective curve. The reason is the Riemann--Roch theorem: the line bundle $L$
with $c_1(L)=\delta$ has $\chi(L)=2+\tfrac12\delta^2=1>0$, so $L$ or its dual has a section
\source{Huy ch.~8 §2.2, footnote~3, ll.~6486--6488; Asp §2.6, ll.~1203--1206}. If $\mathrm{NS}(X)=\Z\delta$,
that curve is a single smooth $\mathbb P^1$ \source{Huy ch.~7, proof of Prop.~5.5,
ll.~6055--6059}.

\begin{theorem}[Kähler cone; Tier~L]\label{thm:k3lattice:kahler}
The Kähler cone of a K3 surface $X$ is
\begin{equation}\label{eq:k3lattice:kahler}
 \mathcal K_X=\{\alpha\in\mathcal C_X\mid(\alpha\cdot C)>0\text{ for all smooth rational
 curves }C\subset X\} .
\end{equation}
In particular, if $\mathrm{Pic}(X)=0$ then $\mathcal K_X=\mathcal C_X$. \TL{}
\source{Huy ch.~8, Thm.~5.2 and Cor.~5.3, ll.~7209--7218}
\end{theorem}
The condition is necessary for an obvious reason: $(J\cdot C)=\int_CJ$ is the area of $C$ in
the Kähler metric, and areas of holomorphic curves are positive. The content of the theorem
is that it is sufficient.

\subsection{Walls, chambers and the Weyl group}
Each root $\delta\in\mathrm{NS}(X)$ defines a \emph{wall}\index{wall} $\delta^\perp\cap
\mathcal C_X$. The walls are locally finite in $\mathcal C_X$, and the connected components
of their complement are called \emph{chambers}\index{chamber}
\source{Huy ch.~8 §2.2, ll.~6494--6519}. By \cref{thm:k3lattice:kahler} the Kähler cone is
one of the chambers. The reflection\index{reflection}
\begin{equation}\label{eq:k3lattice:refl}
 s_\delta(x)=x+(x\cdot\delta)\,\delta,\qquad\delta^2=-2
 \qquad\text{\source{Huy ch.~8 §2.2, ll.~6483--6493}} ,
\end{equation}
is an integral isometry and an involution (\cref{ch:lattices}, where this is proved on the
page and in Lean). It fixes the wall $\delta^\perp$ pointwise, sends $\delta\mapsto-\delta$,
and preserves the positive cone, because $(s_\delta(x)\cdot x)=x^2+(x\cdot\delta)^2>0$ for
$x\in\mathcal C_X$, so that $x$ and $s_\delta(x)$ lie in the same component. Since $\delta$
is of type $(1,1)$, $s_\delta$ fixes $\Omega$: it is a Hodge isometry. The group $W$
generated by all $s_\delta$ is the \emph{Weyl group}\index{Weyl group} of $X$.

\begin{proposition}[Tier~L]\label{prop:k3lattice:weyl}
$W$ acts simply transitively on the set of chambers, and it is generated by the reflections
in the walls of any one chamber. \TL{}
\source{Huy ch.~8, Remark~2.5 and Prop.~2.6, ll.~6542--6563}
\end{proposition}
The source proves this for a lattice of signature $(1,n)$, such as $\mathrm{NS}(X)$ of an
algebraic K3 surface, and treats the Kähler cone of a general K3 surface by the same method
in its §8.5.

\begin{figure}[t]
\centering
\begin{tikzpicture}[scale=1.9]
 \begin{scope}
 \clip (0,0) circle (1);
 \fill[blue!12] (-0.2,-1) -- (0.75,0.7) -- (1,0.3) -- (1,-1) -- cycle;
 \fill[white] (1.2,-0.15) -- (0.2,-1.1) -- (1.2,-1.1) -- cycle;
 \draw[thick] (-0.2,-1) -- (0.85,0.88);
 \draw[thick] (1.2,-0.15) -- (0.2,-1.1);
 \draw (-1,0.35) -- (0.6,1);
 \draw (-1,-0.3) -- (0.1,-1);
 \draw (-0.75,1) -- (-0.45,-1);
 \draw (0.55,1) -- (1,0.55);
 \end{scope}
 \draw[thick] (0,0) circle (1);
 \node at (0.62,-0.18) {$\mathcal K_X$};
 \node[font=\small] at (-0.02,0.25) {$s_\delta\mathcal K_X$};
 \node[font=\small,rotate=61] at (0.22,-0.43) {$\delta^\perp$};
 \draw[<->,>=Stealth] (0.45,0.42) to[bend right=35] (0.12,0.58);
 \node[font=\small] at (1.55,0.85) {$\partial\mathcal C_X$: $\alpha^2=0$};
 \draw[->,>=Stealth] (1.35,0.72) -- (0.93,0.45);
\end{tikzpicture}
\caption{A section $\alpha^2=1$ of the positive cone $\mathcal C_X\subset\R^{1,19}$, drawn
as a disc (a hyperbolic space in the projective model, where hyperplanes are chords). The
walls $\delta^\perp$ of roots cut it into chambers; the shaded chamber is the Kähler cone,
and the reflection $s_\delta$ maps it to the adjacent chamber. Schematic: for $\rho=20$ there
are in general infinitely many walls, accumulating only at the boundary.}
\label{fig:k3lattice:chambers}
\end{figure}

\Cref{fig:k3lattice:chambers} shows the picture. It completes the Torelli theorem: given a
Hodge isometry $\psi$ between $X$ and $X'$, first replace $\psi$ by $-\psi$ if necessary so
that positive cones correspond; then $\psi(\mathcal K_X)$ is one of the chambers of $X'$, and
by \cref{prop:k3lattice:weyl} exactly one $w\in W$ moves it onto $\mathcal K_{X'}$. The
corrected isometry $w\circ\psi$ satisfies the hypothesis of the second half of
\cref{thm:k3lattice:torelli} and is induced by a unique isomorphism. So the group of Hodge
isometries of $X$ is $\{\pm1\}\times W$ extended by the automorphisms of $X$.

\begin{physicsbox}[Compare the circle]
In \cref{ch:lattices} the T-duality $n\leftrightarrow w$ of a circle was identified with the
reflection $s_{e-f}$ of $U$ in the root $e-f$. The half-line of radii $R>0$ is cut by one
``wall'', the self-dual radius $R=\sqrt{\ap}$ fixed by the reflection, into two chambers
exchanged by it, and either chamber is a fundamental domain. For K3 surfaces the Kähler cone
plays the role of the fundamental domain $R\ge\sqrt{\ap}$, the walls are the loci where a
sphere shrinks to zero area, and on the wall something becomes massless in both cases
(\cref{sec:k3lattice:ade}). The analogy is structural; it is not a statement that $s_\delta$
for a geometric $(-2)$-curve is a T-duality.
\end{physicsbox}

\paragraph{Reflections as monodromies and as generators.} Reflections also arise from
families. If a family of K3 surfaces over a disc acquires an $A_1$ singularity at the centre,
transporting $H^2$ once around the centre acts as $s_\delta$, where $\delta$ is the class of
the exceptional $(-2)$-curve \source{Huy ch.~6 §5, ll.~5398--5401}. This is how
$O^+(\Lambda)$ entered \eqref{eq:k3lattice:Mc}: by a theorem essentially due to Wall, every
element of $O^+(\Lambda)$ is a product of reflections $s_{\delta_i}$ in roots
\source{Huy ch.~7 §5.4, ll.~6041--6044}, and each reflection is a monodromy, hence induced
by a diffeomorphism. Note the two different uses of roots. For the monodromy statement
\emph{all} roots of the rank-$22$ lattice $\Lambda$ are used. For the Weyl group of a given
surface only the roots in $\mathrm{NS}(X)$, which depend on the period, are used.

%=====================================================================================
\section{The moduli space of Einstein metrics}\label{sec:k3lattice:einstein}

\subsection{From a two-plane to a three-plane}
A Kähler metric on a K3 surface whose Ricci tensor is proportional to the metric is in fact
Ricci-flat, so ``Einstein'' and ``Ricci-flat'' can be used interchangeably
\source{Asp §2.4, ll.~787--789}. By Yau's theorem\index{Yau's theorem}, for a fixed complex
structure and a fixed Kähler class there is exactly one Ricci-flat Kähler metric
\source{Asp §2.4, ll.~841--844}. The data of such a metric are therefore the two-plane of
$\Omega$ together with a vector $J$ that is
\begin{itemize}[leftmargin=*,itemsep=1pt]
\item orthogonal to the plane, because $J$ is of type $(1,1)$ and $J\wedge\Omega$ is of type
 $(3,1)$;
\item space-like, because $J^2=\int_XJ\wedge J$ is, up to normalization, the volume of $X$
 \source{Asp §2.4, eq.~(38), ll.~849--856}.
\end{itemize}
Thus a Ricci-flat Kähler metric determines a positive \emph{three}-plane
\begin{equation}\label{eq:k3lattice:Sigma}
 \Sigma=\langle\operatorname{Re}\Omega,\operatorname{Im}\Omega,J\rangle\subset
 H^2(X,\R)\cong\R^{3,19} ,
\end{equation}
and a number, the volume. Since $b_+=3$, the plane $\Sigma$ is a \emph{maximal} positive
subspace. It has an intrinsic meaning. The Hodge star $*$ of a Riemannian metric on a
four-manifold squares to $+1$ on two-forms, and harmonic two-forms split into self-dual and
anti-self-dual ones: $H^2(X,\R)=H^+\oplus H^-$. From $\alpha\wedge{*\alpha}=|\alpha|^2\,
\mathrm{vol}$ one reads off that $\alpha^2>0$ for $\alpha=*\alpha$ and $\alpha^2<0$ for
$\alpha=-{*\alpha}$, so $\dim H^+=3$ and $\dim H^-=19$
\source{Asp §2.4, eqs.~(33)--(36), ll.~790--817}. In a local orthonormal coframe
$e^1,\dots,e^4$ with $\dd z_1=e^1+\ii e^2$, $\dd z_2=e^3+\ii e^4$ the self-dual forms are
spanned by
\begin{equation}\label{eq:k3lattice:sd}
 s_1=e^{12}+e^{34},\qquad s_2=e^{13}+e^{42},\qquad s_3=e^{14}+e^{23},
 \qquad e^{ij}:=e^i\wedge e^j ,
\end{equation}
and one checks directly that $\tfrac{\ii}{2}(\dd z_1\wedge\dd\bar z_1+\dd z_2\wedge\dd\bar
z_2)=s_1$ is the Kähler form, while $\dd z_1\wedge\dd z_2=(e^{13}-e^{24})+\ii(e^{14}+e^{23})
=s_2+\ii s_3$ is the holomorphic two-form. Hence
\begin{equation}\label{eq:k3lattice:SigmaHplus}
 \Sigma=H^+\qquad\text{\source{Asp §2.4, eq.~(37) and ll.~825--840}} .
\end{equation}

\begin{figure}[!ht]
\centering
\begin{tikzpicture}[scale=1.5,>=Stealth]
 \draw[thick] (0,0) circle (1);
 \draw[dashed] (-1,0) arc (180:0:1 and 0.3);
 \draw (-1,0) arc (180:360:1 and 0.3);
 \draw[->,thick] (0,0) -- (0,1.45) node[above] {$J$};
 \draw[->,thick] (0,0) -- (1.45,-0.12) node[right] {$\operatorname{Im}\Omega$};
 \draw[->,thick] (0,0) -- (-0.75,-0.62) node[below left] {$\operatorname{Re}\Omega$};
 \draw[->,thick,blue!60!black] (0,0) -- (0.62,0.78);
 \fill[blue!60!black] (0.62,0.78) circle (0.035) node[above right] {$J'$};
 \node at (-1.45,1.0) {$\Sigma=H^+\cong\R^{3,0}$};
 \node[draw,rounded corners,align=center,font=\small] at (3.9,0.35)
  {$\Sigma^\perp=H^-\cong\R^{0,19}$\\[2pt] roots $\delta\in\Gamma^{3,19}$ with
  $\delta\perp\Sigma$\\ $\Longleftrightarrow$ shrunken spheres};
\end{tikzpicture}
\caption{The positive three-plane $\Sigma\subset\R^{3,19}$ of a Ricci-flat metric. A point
$J'$ of the unit sphere in $\Sigma$ selects a Kähler direction, and the plane orthogonal to
it inside $\Sigma$ is the period plane of the corresponding complex structure. All points of
the sphere give the same metric.}
\label{fig:k3lattice:sigma}
\end{figure}

\begin{physicsbox}[One metric, a two-sphere of complex structures]
\index{hyperkähler structure}
The metric determines $\Sigma=H^+$ but \emph{not} its splitting into $\Omega$ and $J$.
Any unit vector $\hat\jmath\in\Sigma$ may be declared to be the direction of the Kähler form;
the orthogonal plane in $\Sigma$ then plays the role of $\Omega$. A K3 surface regarded as a
Riemannian manifold therefore carries a whole two-sphere $S^2$ of complex structures, all
with the same Ricci-flat metric \source{Asp §2.4, ll.~864--874}. This is the hyperkähler
structure: the holonomy $SU(2)\cong Sp(1)$ of \cref{ch:k3} commutes with three complex
structures $I_1,I_2,I_3$ obeying the quaternion relations, and the combinations
$aI_1+bI_2+cI_3$ with $a^2+b^2+c^2=1$ form the $S^2$. For the string this means that the splitting into ``complex structure moduli''
and ``Kähler moduli'', which organizes Calabi--Yau threefolds, has no invariant meaning on
K3. \Cref{fig:k3lattice:sigma} illustrates the geometry.
\end{physicsbox}

\subsection{The moduli space}
\begin{theorem}[Moduli of Einstein metrics; Tier~L]\label{thm:k3lattice:einstein}
\index{moduli space!of Einstein metrics on K3}
The moduli space of Einstein metrics on a K3 surface, including orbifold points, is the
Grassmannian of oriented positive three-planes in $\R^{3,19}$, times the volume, modulo the
diffeomorphisms acting on the lattice:
\begin{equation}\label{eq:k3lattice:ME}
\begin{aligned}
 \mathcal M_E&\cong O^+(\Gamma^{3,19})\big\backslash O^+(3,19)\big/\bigl(SO(3)\times
 O(19)\bigr)\times\R_+\\
 &\cong O(\Gamma^{3,19})\big\backslash O(3,19)\big/\bigl(O(3)\times O(19)\bigr)\times\R_+ .
\end{aligned}
\end{equation}
\TL{} \source{Asp §2.4, Theorem~3 and eqs.~(39)--(40), ll.~877--903}
\end{theorem}
The second form follows from the first because the extra generator $-\mathbf1$ is central
in $O(3,19)$, so adjoining it on the left and on the right changes nothing
\source{Asp §2.4, ll.~900--903}. Let us count dimensions:
\begin{equation}\label{eq:k3lattice:dim58}
 \dim\mathcal M_E=\dim O(3,19)-\dim O(3)-\dim O(19)+1=231-3-171+1=58 .
\end{equation}
Two independent checks. (i) A positive three-plane near a given $\Sigma$ is the graph of a
linear map $\Sigma\to\Sigma^\perp$, which has $3\cdot19=57$ entries; add $1$ for the
volume. (ii) Complex structures have $40$ real parameters \eqref{eq:k3lattice:dim40} and
Kähler classes $20$; the pairs $(\Omega,J)$ fibre over the metrics with fibre $S^2$, and
$40+20-2=58$. Both agree with \eqref{eq:k3lattice:dim58}.

In contrast to \eqref{eq:k3lattice:Mc}, the stabilizer $O(3)\times O(19)$ is now compact, the
discrete group acts properly discontinuously, and $\mathcal M_E$ is Hausdorff
\source{Asp §2.4, ll.~903--906}. Adding the Kähler class has repaired the defect of
$\mathcal M_c$. Comparing with \cref{ch:narain}, $\mathcal M_E$ without the volume factor has
exactly the form of a Narain moduli space with $(d,d)$ replaced by $(3,19)$: a maximal
positive subspace moving relative to an even self-dual lattice. In the torus case the
positive subspace was the space of left-moving momenta; here it is the space of self-dual
harmonic forms. \Cref{ch:stringsk3} adds the $B$-field and the volume to $\Sigma$, obtaining
a positive four-plane in $\R^{4,20}$ and a moduli space of real dimension $4\cdot20=80$ that
is literally of Narain type.

%=====================================================================================
\section{ADE singularities and enhanced gauge symmetry}\label{sec:k3lattice:ade}

\subsection{Where the metric degenerates}
The theorems above included ``orbifold points''. We can now say where they are. Let
$\delta\in\Gamma^{3,19}$ be a root with $\delta\perp\Sigma$. Then $\delta\perp\Omega$, so
$\delta\in\mathrm{NS}(X)$ by \eqref{eq:k3lattice:NS}, and $\pm\delta$ is the class of a
rational curve $C$. But also $\delta\perp J$, so the area $(J\cdot C)$ vanishes: by
\cref{thm:k3lattice:kahler} $J$ is not a Kähler class but lies on a wall of the Kähler cone,
and the sphere $C$ has shrunk to a point.

\begin{theorem}[Orbifold points; Tier~L]\label{thm:k3lattice:orbifold}\index{orbifold}
A point of $\mathcal M_E$ corresponds to an orbifold if and only if the three-plane $\Sigma$
is orthogonal to a root of $\Gamma^{3,19}$. \TL{}
\source{Asp §2.6, Theorem~4 and its proof, ll.~1203--1217}
\end{theorem}
Curves of genus $g\ge1$ cannot shrink in this way: their classes have $C^2=2g-2\ge0$ and
cannot be orthogonal to $\Sigma$, whose orthogonal complement is negative definite
\source{Asp §2.6, ll.~1214--1216}. The simplest local model is the blow-down of one
$(-2)$-curve, which produces the singularity $\C^2/\Z_2$, that is $x_0x_1=x_2^2$ in $\C^3$
with $x_0=\xi^2,x_1=\eta^2,x_2=\xi\eta$ \source{Asp §2.6, eq.~(50), ll.~1082--1090}. The
process is continuous in the moduli space: keep the complex structure fixed and move $J$
until $J\cdot C=0$ \source{Asp §2.6, ll.~1196--1202}.

In general, let $R_\Sigma:=\{\delta\in\Gamma^{3,19}\mid\delta^2=-2,\ \delta\perp\Sigma\}$.
These roots lie in the negative definite space $\Sigma^\perp$, so they form a finite root
system, and the Gram matrix of a set of simple roots has $-2$ on the diagonal and is negative
definite. This is, up to sign, the defining property of the Cartan matrix of a simply-laced
Lie algebra, and the classification is the same: connected components are of type $A_n$,
$D_n$, $E_6$, $E_7$, $E_8$\index{ADE classification} \source{Asp §2.6, ll.~1218--1228}.
Geometrically, the simple roots of one component are spheres which intersect according to
the Dynkin diagram and which are all contracted to one point. For $A_{n-1}$ the singularity
is $\C^2/\Z_n$, that is $x_0x_1=x_2^n$, and its resolution is a chain of $n-1$ spheres
\source{Asp §2.6, ll.~1229--1330}; see \cref{fig:k3lattice:chain}.

\begin{figure}[t]
\centering
\begin{tikzpicture}[scale=0.95,>=Stealth]
 \foreach \k/\a in {0/35,1/-35,2/35} {
   \draw[thick,rotate around={\a:(1.1*\k,0)}] (1.1*\k-0.95,0) -- (1.1*\k+0.95,0);
 }
 \node[font=\small] at (-0.75,-0.8) {$E_1$};
 \node[font=\small] at (1.1,0.85) {$E_2$};
 \node[font=\small] at (2.95,0.85) {$E_3$};
 \draw[->] (4.1,0) -- (5.1,0);
 \foreach \k in {0,1,2} { \fill (5.9+0.9*\k,0) circle (0.09); }
 \draw (5.9,0) -- (7.7,0);
 \node[font=\small] at (6.8,-0.45) {$A_3$};
 \node at (10.2,0) {$-\begin{pmatrix}2&-1&0\\-1&2&-1\\0&-1&2\end{pmatrix}$};
\end{tikzpicture}
\caption{The resolution of $\C^2/\Z_4$: three $(-2)$-curves (drawn as lines), adjacent ones
meeting transversally in one point; the dual graph, which is the Dynkin diagram $A_3$; and
the intersection matrix $E_i\cdot E_j$, minus the Cartan matrix of $SU(4)$.}
\label{fig:k3lattice:chain}
\end{figure}

\begin{example}[A maximally singular K3 surface]\label{ex:k3lattice:maximal}
Continue \cref{ex:k3lattice:rho20}: $\Omega=(e_1+f_1)+\ii(e_2+f_2)$, and choose $J=e_3+f_3$.
Then $J^2=2>0$ and $J\perp\Omega$, so $\Sigma=\langle e_1+f_1,e_2+f_2,e_3+f_3\rangle$ is a
positive three-plane with Gram matrix $2\cdot\mathbf 1_3$, and
\[
 \Sigma^\perp\cap\Gamma^{3,19}=\Z(e_1-f_1)\oplus\Z(e_2-f_2)\oplus\Z(e_3-f_3)\oplus
 E_8(-1)^{\oplus2}\;=\;A_1(-1)^{\oplus3}\oplus E_8(-1)^{\oplus2},
\]
a negative definite lattice of rank $19$ and discriminant $(-2)^3=-8$. In an orthogonal sum
of negative definite even lattices a vector of square $-2$ has all its components but one
equal to zero (the squares of the components are even, non-positive, and add up to $-2$).
So $R_\Sigma$ is the union of the roots of the summands: $2+2+2+240+240=486$ roots, a root
system of type $3A_1+2E_8$ and rank $19$, the largest rank that $\Sigma^\perp\cong\R^{0,19}$
allows. The corresponding point of $\mathcal M_E$ is a K3 orbifold with two $E_8$
singularities and three $\C^2/\Z_2$ singularities, in which $19$ independent spheres have
been shrunk. Moving $J$ to $J_t=e_3+tf_3$ with $t\ne1$, $t>0$, gives
$J_t\cdot(e_3-f_3)=t-1\ne0$ and resolves one of the $A_1$ points, the sphere acquiring area
$|t-1|$. (Symbolic check of the rank, the discriminant and the orthogonality relations:
sympy. The count $240$ is Tier~L, \cref{ch:lattices}.)
\end{example}

\subsection{What the string sees}
Compactify the type~IIA string on a K3 surface (\cref{ch:stringsk3}). The Ramond--Ramond
three-form, reduced on the $22$ two-cycles, gives $22$ abelian gauge fields; the one-form
and the dual of the three-form give two more, so the gauge group is generically $U(1)^{24}$,
and perturbative string theory cannot produce charged states for it
\source{Asp §4.3, ll.~2433--2442}. The heterotic string on $T^4$ has the same moduli space
and also $U(1)^{24}$, but there non-abelian gauge groups appear by the mechanism of
\cref{ch:tduality}: a lattice vector $\alpha$ with $\alpha^2=-2$ becomes orthogonal to the
positive plane, the state has $p_L=0$, $p_R^2=-2$ in the conventions of the source, and it is
a massless vector boson \source{Asp §4.3, ll.~2450--2483}. If the two theories are dual, the
type~IIA string must do the same.

\begin{proposition}[Aspinwall's Proposition~2; Tier~L, conditional]
\label{prop:k3lattice:gauge}\index{enhanced gauge symmetry}
If a K3 surface has an orbifold singularity, a type~IIA string compactified on it can
exhibit a non-abelian gauge group, and the ADE type of the singularity is the ADE type of the
gauge algebra; for instance $\C^2/\Z_n$ gives $SU(n)$. \TL{}
\source{Asp §4.3, Proposition~2, ll.~2525--2531}
\end{proposition}
Two qualifications belong to the statement. First, it rests on the duality between the
type~IIA string on K3 and the heterotic string on $T^4$, which the source presents as a
proposition that cannot at present be proved, together with the assumptions made in
constructing the moduli spaces \source{Asp §4.3, ll.~2527--2528}. Its tier is that of a
well-tested conjecture of the literature. Secondly, the singularity is necessary but not
sufficient: the component of the $B$-field along the vanishing cycle must also vanish, while
the solvable orbifold conformal field theories have $B=\tfrac12$ there. This is why string
theory on an orbifold is perfectly regular and yet shows no enhanced symmetry
\source{Asp §4.3, ll.~2530--2546}. In \cref{ex:k3lattice:maximal} the enhanced gauge algebra
at $B=0$ would be $E_8\oplus E_8\oplus SU(2)^3$, of rank $19$, inside the total rank $24$.

\begin{pitfallbox}[Three meanings of ``$E_8$'' in this chapter]
(i) The summand $E_8(-1)$ of the abstract lattice $\Lambda$, present for every K3 surface,
including smooth ones with no curves at all. (ii) An $E_8$ \emph{singularity}, present only
at special points of $\mathcal M_E$, where eight spheres with the $E_8$ intersection matrix
have zero area. (iii) An $E_8$ \emph{gauge symmetry} of the type~IIA string, present only if
in addition the $B$-field vanishes on these spheres. The implication is (iii) $\Rightarrow$
(ii) $\Rightarrow$ (i) and none of the converses holds. The Lean results below concern (i)
only.
\end{pitfallbox}

%=====================================================================================
\section{What the machine has checked}\label{sec:k3lattice:lean}

Almost everything in this chapter is geometry and analysis on complex manifolds, and none
of that is formalized: the Lean libraries of this book contain no manifolds, no cohomology
and no Hodge structures. What is formalized is the finite arithmetic which the geometry rests
on. Three files of the directory \lean{DualScaleStream2/Lattice/} are relevant:
\lean{K3T2Signature.lean}, \lean{E8.lean} and \lean{Reflection.lean}. The last two were
presented in \cref{ch:lattices}; here we describe what they contribute to the K3 lattice. All
theorems quoted depend only on the axioms \lean{propext}, \lean{Classical.choice} and
\lean{Quot.sound}.

\begin{leanbox}[Signature bookkeeping (\lean{K3T2Signature.lean})]
\begin{leancode}
structure Signature where
  pos : ℕ
  neg : ℕ
  deriving DecidableEq, Repr

namespace Signature
instance : Add Signature := ⟨fun a b => ⟨a.pos + b.pos, a.neg + b.neg⟩⟩
def rank (s : Signature) : ℕ := s.pos + s.neg
def index (s : Signature) : ℤ := (s.pos : ℤ) - s.neg
end Signature

def sigU : Signature := ⟨1, 1⟩
def sigE8Neg : Signature := ⟨0, 8⟩
def sigK3 : Signature := sigE8Neg + sigE8Neg + sigU + sigU + sigU
def sigMukai : Signature := sigK3 + sigU

theorem sigK3_eq : sigK3 = ⟨3, 19⟩ := by rfl
theorem rank_K3 : sigK3.rank = 22 := by ...
theorem index_mod_eight :
    sigK3.index % 8 = 0 ∧ sigMukai.index % 8 = 0 ∧ sigK3T2.index % 8 = 0 := by ...
\end{leancode}
\textbf{Reading.} A \lean{Signature} is a bare pair of natural numbers. Nothing in the
structure ties it to a quadratic form. \lean{sigU} and \lean{sigE8Neg} are
\emph{definitions}: the numbers $(1,1)$ and $(0,8)$ are entered by hand. They are justified
on paper by Sylvester's law together with the Lean theorems \lean{hyperbolicU\_congruence}
and \lean{cartanE8\_posDef} of \cref{ch:lattices}, but no Lean statement derives the pairs
from those theorems, since the library has no definition of the signature of a matrix.
\lean{sigK3\_eq} then says: adding the pair $(0,8)$ twice and the pair $(1,1)$ three times
gives $(3,19)$. \lean{rank\_K3} says $3+19=22$. \lean{index\_mod\_eight} says that
$3-19=-16$, $4-20=-16$ and $6-22=-16$ are divisible by $8$ (for integers Lean's \lean{\%}
returns a non-negative remainder, and $-16=(-2)\cdot8+0$); \lean{sigK3T2}, of value $(6,22)$,
belongs to \cref{ch:k3t2}. This verifies the necessary condition of the classification
theorem at three points. It is not the theorem that every even unimodular lattice satisfies
$n_+\equiv n_-\pmod8$, which remains Tier~L.

\textbf{Proof idea.} Each statement reduces by unfolding definitions to a closed equation
between numerals, which the kernel evaluates: \lean{rfl} for \lean{sigK3\_eq} and for the
three conjuncts of \lean{index\_mod\_eight}, \lean{simp} with the definitions followed by
\lean{rfl} for \lean{rank\_K3}.
\end{leanbox}

\begin{leanbox}[Two routes to $(3,19)$ compared (\lean{K3T2Signature.lean})]
\begin{leancode}
theorem sigK3_matches_hodge :
    sigK3.pos = StringTheory.UseCases.K3Signature.bPlus ∧
    sigK3.neg = StringTheory.UseCases.K3Signature.bMinus := by ...
\end{leancode}
\textbf{Reading.} In an older library of the project the numbers \lean{bPlus} and
\lean{bMinus} are defined from a table of the Hodge numbers of K3 by the formulas
\eqref{eq:k3lattice:bpm}, $b_+=2h^{2,0}+1$ and $b_-=h^{1,1}-1$, and the theorems
\lean{bPlus\_eq\_three}, \lean{bMinus\_eq\_nineteen} and
\lean{k3\_signature\_eq\_neg\_sixteen} evaluate them. The present theorem says that the
lattice-decomposition route and the Hodge route produce the same two natural numbers. This
is the formal counterpart of the paragraph ``Signature $(3,19)$, by two routes'' in
\cref{sec:k3lattice:form}. Both inputs are tabulated Tier~L facts. Their agreement is a
consistency check of the bookkeeping; it is not a derivation of the signature of a K3
surface, and the formulas \eqref{eq:k3lattice:bpm} themselves (the Hodge index theorem) are
not proved in Lean.

\textbf{Proof idea.} Both sides unfold to the numerals $3$ and $19$; \lean{rfl} closes each
half.
\end{leanbox}

\begin{leanbox}[The summand $E_8(-1)$ (\lean{E8.lean})]
\begin{leancode}
def e8Neg : Gram 8 := -cartanE8

theorem e8Neg_symm : e8Negᵀ = e8Neg := by ...
theorem e8Neg_evenDiag : IsEvenDiag e8Neg := by ...
theorem e8Neg_unimodular : IsUnimodular e8Neg := by ...
\end{leancode}
\textbf{Reading.} \lean{Gram 8} is the type of $8\times8$ integer matrices and
\lean{cartanE8} is the Cartan matrix of \cref{ch:lattices}. The three theorems say that the
matrix $-E_8$ is symmetric, has even diagonal entries (all equal to $-2$), and has
determinant $\pm1$. Combined with the general lemma
\lean{even\_quadratic\_form\_of\_even\_diag} of \cref{ch:lattices} the first two give that
$x^{\mathsf T}(-E_8)x$ is even for every integer vector $x$. These are exactly the three
properties that \cref{sec:k3lattice:form} derived for $H^2(X,\Z)$ from Poincaré duality and
Wu's formula, verified here for one block of the right-hand side of
\eqref{eq:k3lattice:lattice}. The same properties of $U$ are \lean{hyperbolicU\_symm},
\lean{hyperbolicU\_evenDiag} and \lean{hyperbolicU\_unimodular}. The $22\times22$ matrix
$G_\Lambda$ of \cref{ex:k3lattice:gram} is not assembled in Lean, and no Lean statement
asserts that a direct sum of even unimodular Gram matrices is even unimodular, although on
paper this is immediate: the diagonal is the union of the diagonals and the determinant is
the product.

\textbf{Proof idea.} Symmetry and the diagonal are decided entry by entry (\lean{fin\_cases}).
Unimodularity uses a certificate: the integer matrix \lean{cartanE8Inv} satisfies
\lean{cartanE8Inv * cartanE8 = 1} (theorem \lean{cartanE8Inv\_mul}, $64$ entries checked by
the kernel), hence $(-E_8^{-1})(-E_8)=\mathbf1$, and \lean{isUnimodular\_of\_mul\_eq\_one}
converts an integer inverse into $\det=\pm1$. Note that this is the argument used for
\eqref{eq:k3lattice:PD}: Poincaré duality provides an integer inverse of the Gram matrix.
\end{leanbox}

\begin{leanbox}[Reflections in $(-2)$-vectors (\lean{Reflection.lean})]
\begin{leancode}
def reflection (L : Gram n) (v : Fin n → ℤ) : Gram n :=
  1 + vecMulVec v v * L
def latticeNorm (L : Gram n) (v : Fin n → ℤ) : ℤ := v ⬝ᵥ (L *ᵥ v)

theorem reflection_isometry (L : Gram n) (hL : Lᵀ = L) (v : Fin n → ℤ)
    (hv : latticeNorm L v = -2) :
    (reflection L v)ᵀ * L * reflection L v = L := by ...

theorem e8Neg_simpleRoot_norm (i : Fin 8) :
    latticeNorm e8Neg (Pi.single i 1) = -2 := by ...

theorem e8Neg_weyl_isometry (i : Fin 8) :
    (reflection e8Neg (Pi.single i 1))ᵀ * e8Neg
      * reflection e8Neg (Pi.single i 1) = e8Neg := ...
\end{leancode}
\textbf{Reading.} \lean{reflection L v} is the matrix $\mathbf1+vv^{\mathsf T}L$ of the map
\eqref{eq:k3lattice:refl}. \lean{reflection\_isometry} holds for every rank $n$ and every
symmetric integer matrix $L$, so in particular for the $22\times22$ matrix $G_\Lambda$ and
any root $\delta$ of the K3 lattice, even though $G_\Lambda$ is not defined in Lean: the
reflection in a vector of square $-2$ preserves the form. Its companion
\lean{reflection\_involution} gives $s_\delta^2=\mathbf1$. \lean{Pi.single i 1} is the $i$-th
standard basis vector, that is, the $i$-th simple root $\alpha_i$.
\lean{e8Neg\_simpleRoot\_norm} says $\alpha_i^2=-2$, and \lean{e8Neg\_weyl\_isometry} says
that each of the eight simple reflections preserves $-E_8$. The theorem covers eight
reflections. It does not define the Weyl group, does not show that these eight generate it,
and says nothing about the other $232$ roots, though \lean{reflection\_isometry} would apply
to each of them once its norm is computed. No Lean statement connects these matrices with
curves on a surface, with monodromy, or with gauge bosons.

\textbf{Proof idea.} The key identity is $(vv^{\mathsf T}L)(vv^{\mathsf T}L)=
(v^{\mathsf T}Lv)\,vv^{\mathsf T}L$ (\lean{vecMulVec\_mul\_self}); at norm $-2$ the quadratic
term in the expansion of $(\mathbf1+Lvv^{\mathsf T})L(\mathbf1+vv^{\mathsf T}L)$ cancels the
two linear terms. The $E_8$ statements are instances: the norm is computed by
\lean{fin\_cases} and \lean{simp}, and \lean{e8Neg\_weyl\_isometry} is
\lean{reflection\_isometry} applied to \lean{e8Neg\_symm} and \lean{e8Neg\_simpleRoot\_norm}.
\end{leanbox}

\paragraph{What a faithful formalization would need.} To state even
\cref{thm:k3lattice:lattice} in Lean one needs singular cohomology of a manifold with its cup
product, Poincaré duality, characteristic classes for Wu's formula and the signature
theorem, and the classification of indefinite unimodular forms. To state the Torelli theorem
one needs in addition complex manifolds, Hodge decomposition and deformation theory. None of
this was available in the Mathlib version used by the project, and the programme did not
attempt it. A realistic intermediate goal, purely algebraic, would be (compare \cref{ch:open}):
define lattices as $\Z$-modules with a form, define isometries and direct sums, prove that
$G_\Lambda$ is even and unimodular, and prove that $s_\delta\in O(\Lambda)$ generate a group
acting on the set of roots. Exercises~\ref{ex:k3lattice:exlean1} and~\ref{ex:k3lattice:exlean2} are first steps.

\begin{center}
\small
\begin{tabular}{@{}p{0.47\textwidth}cp{0.40\textwidth}@{}}
\toprule
Statement & Tier & Where\\
\midrule
$(0,8)+(0,8)+3\,(1,1)=(3,19)$; rank $22$; $-16\equiv0\ (8)$ & \TA &
 \lean{sigK3\_eq}, \lean{rank\_K3}, \lean{index\_mod\_eight}\\[2pt]
Lattice route and Hodge route give the same pair & \TA & \lean{sigK3\_matches\_hodge}\\[2pt]
$-E_8$ symmetric, even diagonal, $\det=\pm1$ & \TA & \lean{e8Neg\_symm},
 \lean{e8Neg\_evenDiag}, \lean{e8Neg\_unimodular}\\[2pt]
$s_v$ is an isometry and an involution if $v^{\mathsf T}Lv=-2$; the eight simple
 reflections of $-E_8$ & \TA & \lean{reflection\_isometry}, \lean{reflection\_involution},
 \lean{e8Neg\_weyl\_isometry}\\[2pt]
$G_\Lambda$: $\det=-1$, inertia $(3,19)$; \cref{ex:k3lattice:rho20,ex:k3lattice:maximal} &
 --- & symbolic checks (sympy) and derivations on the page\\[2pt]
$H^2(X,\Z)$ is even, unimodular, of signature $(3,19)$, hence $\cong\Lambda$ & \TL &
 Huy ch.~1 Prop.~3.5; Asp §2.3\\[2pt]
Local Torelli, surjectivity, global Torelli; Kähler cone; $W$ simply transitive on
 chambers; $\mathrm{Mon}(X)=O^+$ & \TL & Huy ch.~6--8\\[2pt]
$\mathcal M_c$, $\mathcal M_E$ as double cosets; orbifold points $\Leftrightarrow$
 root $\perp\Sigma$; ADE & \TL & Asp §2.3--2.6\\[2pt]
ADE singularity with $B=0$ $\Rightarrow$ non-abelian gauge symmetry of type IIA & \TL &
 Asp §4.3, conditional on IIA/heterotic duality\\[2pt]
A Lean \lean{Signature} or Gram matrix ``is'' the cohomology of a K3 surface & --- &
 an identification, never Tier~A\\
\bottomrule
\end{tabular}
\end{center}

%=====================================================================================
\begin{summarybox}
\begin{itemize}[leftmargin=*,itemsep=1pt]
\item $H^2(X,\Z)$ of a K3 surface with the intersection pairing is a lattice of rank $22$. It
 is unimodular by Poincaré duality (the dual basis is an integer inverse of the Gram matrix),
 even by Wu's formula and $c_1=0$, and of signature $(3,19)$ by the signature theorem or by
 $b_+=2h^{2,0}+1$. Classification then gives $H^2(X,\Z)\cong3U\oplus2E_8(-1)$, abstractly.
\item The class $\Omega=x+\ii y$ satisfies $\Omega^2=0$, $\Omega\cdot\bar\Omega>0$:
 $x\perp y$, $x^2=y^2>0$. The period is the point $[\Omega]$ of the $20$-dimensional period
 domain $D\cong O(3,19)/(SO(2)\times O(1,19))$, equivalently an oriented positive two-plane.
 It determines the Hodge decomposition and the Picard lattice $\Omega^\perp\cap H^2(X,\Z)$.
\item Torelli theorems (Tier~L): the period map is a local isomorphism, surjective, and a
 Hodge isometry implies isomorphism. A Hodge isometry comes from a unique isomorphism if it
 respects Kähler cones. The moduli space $\mathcal M_c$ of complex structures is a double
 coset by $O^+(\Gamma^{3,19})$ and is not Hausdorff.
\item Roots $\delta^2=-2$ of $\mathrm{NS}(X)$ are, up to sign, classes of curves; $(-2)$-curves
 bound the Kähler cone, which is one chamber of the positive cone; the Weyl group generated
 by $s_\delta(x)=x+(x\cdot\delta)\delta$ permutes the chambers simply transitively.
\item A Ricci-flat metric is a positive three-plane $\Sigma=H^+=\langle\Omega,J\rangle$ and a
 volume: $\mathcal M_E\cong O(\Gamma^{3,19})\backslash O(3,19)/(O(3)\times O(19))\times\R_+$, of
 dimension $57+1=58$, Hausdorff, of Narain type. One metric carries an $S^2$ of complex
 structures.
\item $\Sigma\perp$ root $\Leftrightarrow$ orbifold point; the roots orthogonal to $\Sigma$
 form an ADE root system of rank $\le19$; with vanishing $B$-field the type~IIA string has
 the corresponding non-abelian gauge symmetry (conditional on string duality).
\item Lean checks the signature arithmetic, the evenness and unimodularity of the block
 $-E_8$, and that reflections in $(-2)$-vectors are involutive isometries. It does not
 contain K3 surfaces, periods or the Torelli theorem.
\end{itemize}
\end{summarybox}

%=====================================================================================
\section*{Exercises}
\addcontentsline{toc}{section}{Exercises}

\begin{exercise}[$\star$]\label{ex:k3lattice:exsig}
Solve $b_++b_-=22$, $b_+-b_-=-16$. Then repeat the computation for the four-torus $T^4$,
which has $\chi=0$, $c_1=0$ and $b_2=6$: find its signature, and use the classification of
\cref{ch:lattices} to show that $H^2(T^4,\Z)\cong3U$, given that it is even. Compare with the
K3 lattice: which summand is ``new''?
\end{exercise}

\begin{exercise}[$\star$]
A smooth quartic surface may contain a line $\ell\cong\mathbb P^1$. Check that the Fermat
quartic $x_0^4+x_1^4+x_2^4+x_3^4=0$ contains the line $x_0=\zeta x_1$, $x_2=\zeta x_3$ for
$\zeta^4=-1$. Use \eqref{eq:k3lattice:adjunction} to find
$\ell^2$, and compute the Gram matrix of the sublattice of $\mathrm{NS}(X)$ spanned by the
hyperplane class $h$ and $\ell$, using $h\cdot\ell=1$. Verify that its signature is $(1,1)$,
in accordance with the signature $(1,\rho-1)$ of the Picard lattice.
\end{exercise}

\begin{exercise}[$\star\star$]
Let $\Omega=x+\ii y$ with $x=e_1+2f_1$ and $y=\sqrt2\,(e_2+f_2)$ in the notation of
\cref{ex:k3lattice:rho20}. (a) Verify that $[\Omega]\in D$. (b) Show that a lattice vector
$v=ae_1+bf_1+ce_2+df_2+\dots$ is orthogonal to $\Omega$ if and only if $2a+b=0$ and $c+d=0$,
and conclude that $\rho=20$ again. (c) Replace $\sqrt2$ by a transcendental number $t$ and
$x$ by $e_1+t'f_1$ with $t'$ chosen such that $x^2=y^2$. What is $\rho$ now? The Picard number is not constant on $D$: it jumps at special periods.
\end{exercise}

\begin{exercise}[$\star\star$]
Let $\delta,\delta'$ be roots of a lattice and $g$ an isometry. (a) Prove
$g\circ s_\delta=s_{g(\delta)}\circ g$. (b) If $\delta\cdot\delta'=0$, show that $s_\delta$
and $s_{\delta'}$ commute. (c) If $\delta\cdot\delta'=1$, show that $s_\delta(\delta')
=\delta+\delta'$ is a root and that $(s_\delta s_{\delta'})^3=1$. These are the relations of
the Weyl group of $A_2$, the symmetric group $S_3$. Which singularity and which gauge group
correspond to this configuration?
\end{exercise}

\begin{exercise}[$\star\star$]
In \cref{ex:k3lattice:maximal} replace $J$ by $J'=e_3+f_3+\epsilon\,w$, where $w$ is a vector
in the first $E_8(-1)\otimes\R$ with $(w\cdot\alpha_i)\ne0$ for all simple roots except
$\alpha_0$ and $\alpha_1$, and $\epsilon$ is small. Check that $J'$ is still orthogonal to
$\Omega$ and space-like. Which roots of the first $E_8(-1)$ remain orthogonal to $\Sigma$?
Describe the singularities of the resulting orbifold and the gauge algebra at $B=0$. (Hint:
a root orthogonal to $w$ must be a combination of simple roots orthogonal to $w$, if $w$ is
chosen in the closure of the fundamental chamber.)
\end{exercise}

\begin{exercise}[$\star\star$, Lean]\label{ex:k3lattice:exlean1}
State and prove in Lean, in the style of \lean{sigK3\_eq}, that the signature pair of
$3U$ is $(3,3)$, that its rank is $6$ and that its index is divisible by $8$. Then state
(without proving anything about lattices) the pair for $E_8(-1)^{\oplus3}\oplus U^{\oplus k}$
and find, by \lean{decide}, for which $k\le4$ the rank equals $26$. Which part of
Exercise~\ref{ex:k3lattice:exsig} is invisible to this arithmetic?
\end{exercise}

\begin{exercise}[$\star\star\star$, Lean]\label{ex:k3lattice:exlean2}
Define in Lean the vector \lean{v : Fin 8 → ℤ} of the root $\alpha_3+\alpha_4$ of $E_8(-1)$
and prove \lean{latticeNorm e8Neg v = -2}. Deduce from \lean{reflection\_isometry} that the
reflection in this non-simple root is an isometry of \lean{e8Neg}. Then prove the matrix
identity $s_{\alpha_4}s_{\alpha_3}s_{\alpha_4}=s_{\alpha_3+\alpha_4}$ for the corresponding
\lean{reflection} matrices (a finite computation: \lean{ext i j}, \lean{fin\_cases},
\lean{simp}). This is the identity of part~(a) of the exercise on isometries above with
$g=s_{\alpha_4}$.
\end{exercise}

\begin{exercise}[$\star\star\star$]
(A period domain that is an upper half-plane.) Let $\Lambda'=U\oplus\langle2\rangle$, with
basis $e,f,w$ and $w^2=2$; its signature is $(2,1)$. Using \cref{rem:k3lattice:tube}, show that $D$ is the union of the
upper and lower half-planes with coordinate $\zeta=\tau w$, and that the reflection in the
root $e-f$ acts on the period $[1:-\tau^2:\tau]$ (coordinates along $e,f,w$) by
$\tau\mapsto-1/\tau$. Compare with the T-duality and S-transformations of
\cref{ch:torus2}.
\end{exercise}

%=====================================================================================
\section*{Further reading}
\addcontentsline{toc}{section}{Further reading}
\begin{itemize}[leftmargin=*,itemsep=2pt]
\item The K3 lattice: \source{Huy ch.~1 §3.3, Prop.~3.5 and Remarks~3.6--3.7, ll.~561--660};
 the physicist's derivation with Poincaré duality and the period relations,
 \source{Asp §2.3, ll.~502--774}.
\item Period domains, their Grassmannian and tube-domain descriptions, and the action of
 $O(\Lambda)$: \source{Huy ch.~6 §1, ll.~4272--4420}; local Torelli,
 \source{Huy ch.~6 §2, Prop.~2.8, ll.~4654--4670}.
\item Moduli of marked surfaces, surjectivity of the period map and the global Torelli theorem
 via monodromy: \source{Huy ch.~7 §§2--5, ll.~5594--6060}.
\item Chambers, walls, the Weyl group and the Kähler cone:
 \source{Huy ch.~8 §2, ll.~6423--6563, and §5, ll.~7195--7222}.
\item Einstein metrics, the three-plane $\Sigma$ and hyperkähler rotation:
 \source{Asp §2.4, ll.~777--908}; algebraic K3 surfaces and the transcendental lattice,
 \source{Asp §2.5, ll.~910--1029}; orbifolds, blow-ups and the ADE classification,
 \source{Asp §2.6, ll.~1033--1330}; enhanced gauge symmetry,
 \source{Asp §4.3, ll.~2428--2560}.
\item The three Lean files of \cref{sec:k3lattice:lean} are short and carry detailed module
 documentation.
 The lattice theory behind them is \cref{ch:lattices}; the explicit cycles realizing part of
 $\Lambda$ on a Kummer surface are in \cref{ch:kummer}; the extension to $\Gamma^{4,20}$ is in
 \cref{ch:mukai,ch:stringsk3}.
\end{itemize}
